\documentclass[a4paper,11pt]{texMemo}
\usepackage[english]{babel}
\usepackage{graphicx}
\usepackage{url}

\memoto{Organizing Committee, Atlantic Causal Inference Conference 2019}
\memofrom{Prof.~Mark J.~van der Laan}
\memosubject{Workshop proposal: The \texttt{tlverse} software ecosystem for
  causal inference}
\memodate{\today}
\logo{\includegraphics[scale=0.15]{figs/ucberkeleyseal_874_540.eps}}

\begin{document}
\maketitle

\vspace{-0.25in}
\section{Workshop information}

\begin{itemize}
  \itemsep1pt
  \item \textbf{Title:} ``The \texttt{tlverse}: A Software Ecosystem for Causal
    Inference via Targeted Learning''
  \item \textbf{Goal:} This will primarily be a software workshop centered
    around the new \texttt{tlverse} ecosystem (\url{https://github.com/tlverse})
    of \texttt{R} packages. In order to motivate the use of the software tools,
    there will be substantive discussion of both causal inference methodology
    --- focusing on the field of targeted learning --- and applications in both
    large-scale observational studies and randomized experiments.
\end{itemize}

\section{Abstract}

This full-day short course will provide a comprehensive introduction to both the
\texttt{tlverse} software ecosystem and the field of targeted learning for
causal inference. The workshop will primarily focus on introducing modern
methodological developments in statistical causal inference and their
corresponding software implementations in the \texttt{tlverse}. Topics to be
addressed include efficient estimation in nonparametric and semiparametric
models; efficiency bounds for causal estimands; efficient estimators based on
influence functions, constructed via targeted minimum loss-based estimation;
bias-correction and de-confounding for such estimators via machine learning;
stochastic interventions (e.g., incremental propensity score shifts), and
optimal treatment regimes (and extensions for categorical treatments). The
strategies we discuss yield (double) robust estimators with normal limiting
distributions and fast convergence rates in complex statistical models, allowing
for valid inference even when nuisance parameters are estimated via machine
learning. Each causal parameter and statistical estimator will be demonstrated
with a corresponding \texttt{R} package from the \texttt{tlverse} ecosystem via
hands-on data analysis, providing participants opportunities to familiarize
themselves with both the methods and tools. Some backgound in mathematical
statistics will be useful; familiarity with the \texttt{R} language and
environment for statistical computing will be essential.

\section{Motivations}

Randomized clinical trials (RCTs) have long served as the gold standard of
evidence for comparing potential interventions in clinical medicine, public
health, and marketing, political science, and many other fields. Unfortunately,
such trials are often not feasible due to ethical, logistic or economical
constraints. Observational studies constitute a potentially rich alternative to
RCTs, providing an opportunity to learn about the causal effects of
interventions for which little or no trial data can be produced; however, in
such studies intervention allocation may be strongly confounded by other
important characteristics. Thus, great care is needed in attempts to disentangle
observed relationships and, ultimately, infer causal effects. This course will
provide an overview of recent advances in the field of targeted learning, a
modern statistical framework for that advocates the use of state-of-the-art
machine learning to flexibly adjust for confounding while yielding valid
statistical inference, thus unlocking observational studies for causal
inference.

A longstanding problem in the use of modern, complex statistical approaches
has been the availability of robust software to facilitate computationally
reproducible causal inference analyses. Thus, to complement the statisitcal
methodology introduced, this workshop will focus on the \texttt{tlverse}
software ecosystem, a recent effort to develop a  suite of \texttt{R} packages
that share a core set of design principles centered on extensibility. Tools to
be introduced include the Super Learner model stacking framework to flexibly
adjust for confounding, and its implementation in the \texttt{sl3 R} package;
estimation and inference for ``classical'' causal parameters (e.g., ATE) with
the \texttt{tmle3 R} package; optimal treatment regimes and the
\texttt{tmle3mopttx R} package; stochastic interventions and the
\texttt{tmle3shift R} package; and variable importance analyses via targeted
learning.

\section{Organizers}

\subsection{Mark van der Laan, Ph.D.}

Mark van der Laan, PhD, is Professor of Biostatistics and Statistics at UC
Berkeley. His research interests include statistical methods in computational
biology, survival analysis, censored data, targeted maximum likelihood
estimation, causal inference, data-adaptive loss-based learning, and multiple
testing. His research group developed loss-based super learning in
semiparametric models, based on cross-validation, as a generic optimal tool for
the estimation of infinite-dimensional parameters, such as nonparametric density
estimation and prediction with both censored and uncensored data. Building on
this work, his research group developed targeted maximum likelihood estimation
for a target parameter of the data-generating distribution in arbitrary
semiparametric and nonparametric models, as a generic optimal methodology for
statistical and causal inference. Most recently, Mark's group has focused in
part on the development of a centralized, principled set of software tools for
targeted learning, the \texttt{tlverse}. Contact: \texttt{laan@berkeley.edu}.

\subsection{Alan Hubbard, Ph.D.}

Alan Hubbard is Professor of Biostatistics, former head of the Division of
Biostatistics at UC Berkeley, and head of data analytics core at UC Berkeley's
SuperFund research program. His current research interests include causal
inference, variable importance analysis, statistical machine learning,
estimation of and inference for data-adaptive statistical target parameters, and
targeted minimum loss-based estimation. Research in his group is generally
motivated by applications to problems in computational biology, epidemiology,
and precision medicine. Contact: \texttt{hubbard@berkeley.edu}.

\subsection{Jeremy Coyle, Ph.D.}

Jeremy Coyle is a consulting data scientist and statistical programmer,
currently leading the software development effort that has produced the
\texttt{tlverse} ecosystem of \texttt{R} packages and related software tools.
Jeremy earned his PhD in Biostatistics from UC Berkeley in 2016, primarily under
the supervision of Alan Hubbard. Contact: \texttt{jeremyrcoyle@gmail.com}.

\subsection{Nima Hejazi, M.A.}

Nima is a PhD candidate in biostatistics with a designated emphasis in
computational and genomic biology at UC Berkeley, under the joint supervision of
Mark van der Laan and Alan Hubbard. Nima is affiliated with UC Berkeley's Center
for Computational Biology and NIH Biomedical Big Data training program. Nima's
research interests span causal inference, nonparametric inference and machine
learning, targeted loss-based estimation, survival analysis and censored
data models, statistical computing, reproducible research, and computational
biology. Substantive applications have recently included vaccine efficacy
trials, precision medicine, and high-dimensional biology. Nima is also
passionate about software development for applied statistics, including software
design, automated testing, and reproducible coding practices. Contact:
\texttt{nhejazi@berkeley.edu}.

\subsection{Ivana Malenica, M.A.}

Ivana is a PhD student in the UC Berkeley Biostatistics Division, working
with Alan Hubbard and Mark van der Laan. She earned her Master's in
Biostatistics and Bachelor's in Mathematics, and spent some time at the
Translational Genomics Research Institute. Some of her prior work includes
mathematical modeling, Bayesian models for allele specific expression and
time-series models for genomics data. Broadly, her research interests span
Machine Learning, Causal Inference, high-dimensional data, and semiparametric
theory. Some of her current work has been centered around active learning, model
selection criterion for dependent data, targeted estimators for parameters of
semi- and nonparametric models (recently, working with the natural mediation
effect) and software development (\texttt{medltmle}, \texttt{sl3}). Contact:
\texttt{imalenica@berkeley.edu}.

\section{Duration}

This will be a 6-hour, full-day workshop, featuring modules that each introduce
a distinct causal question, alongside statistical methodology and software for
implementing solutions to the given problem. A sample schedule would take the
form:
\begin{itemize}
  \itemsep0pt
  \item 09:30AM--10:20AM: Introduction to targeted learning for causal inference
  \item 10:20AM--10:30AM: Coffee Break
  \item 10:30AM--11:20AM: Introduction to the \texttt{tlverse} software
    ecosystem
  \item 11:20AM--11:30AM: Coffee Break
  \item 11:30AM--12:45PM: Stochastic interventions and the \texttt{tmle3shift R}
    package
  \item 12:45PM--01:30PM: Lunch Break
  \item 01:30PM--02:45PM: Optimal treatments regimes and the
    \texttt{tmle3mopttx R} package
  \item 02:45PM--02:55PM: Coffee Break
  \item 02:55PM--03:45PM: Targeted learning for variable importance analyses
  \item 03:45PM--4:00PM: Course summary and concluding remarks
\end{itemize}

\section{Prior History}

The \texttt{tlverse} ecosystem is a relatively recent effort (about 2 years in
the making) in developing a set of software tools for causal inference, all
built around a consistent set of design principles. Although we have introduced
the material in graduate courses taught at UC Berkeley, this will potentially be
the first offering in the 6-hour format.

\end{document}

How to add new estimators?